\chapter{Making measurements projective}
\label{chap:projective}


\begin{theorem}[Naimark dilation]\label{thm:naimark}
  \lean{MIPStarRE.Paper2009LDT.Section5MakingMeasurementsProjective.naimark}
  \uses{def:submeasurement}
  Let $\ket{\psi} \in \mathcal H_{\mathrm A} \ot \mathcal H_{\mathrm B}$, and let $A=\{A_a^x\}$ and $B=\{B_b^y\}$ be submeasurements on the two sides. Then there are auxiliary Hilbert spaces, a product auxiliary state $\ket{\mathrm{aux}}$, and projective measurements $\widehat A$ and $\widehat B$ on the enlarged spaces such that, writing $\ket{\widehat\psi}=\ket{\psi}\ot\ket{\mathrm{aux}}$,
  \[
    \bra{\psi} A_a^x \ot B_b^y \ket{\psi} = \bra{\widehat\psi} \widehat A_a^x \ot \widehat B_b^y \ket{\widehat\psi}
  \]
  for every $x,y,a,b$.
\end{theorem}
\begin{proof}
  Apply one-measurement Naimark dilation separately to every family $A^x$ and $B^y$. For each question $x$ on Alice's side and $y$ on Bob's side, introduce an auxiliary register carrying the local dilation of that particular submeasurement, together with one extra outcome for the missing mass. Let $\ket{\mathrm{aux}}$ be the tensor product of all these auxiliary states. The dilated operator $\widehat A_a^x$ acts as the Naimark projector for $A^x$ on the $x$-th auxiliary register and as the identity on the other auxiliary factors, and similarly for $\widehat B_b^y$. Because the two sides use disjoint tensor factors, the correlation for any fixed pair $(x,y)$ reduces to the one-measurement Naimark identity on those two registers, which gives the stated equality.
\end{proof}

\begin{lemma}[Orthogonalization lemma for measurements]\label{lem:orthonormalization-main-lemma}
  \lean{MIPStarRE.Paper2009LDT.Section5MakingMeasurementsProjective.orthonormalizationMainLemma}
  \uses{def:simeq, def:approx_delta}
  Let $A=\{A_a\}$ and $B=\{B_a\}$ be measurements on a state $\ket{\psi}$ such that
  \[
    A_a \ot I \simeq_\zeta I \ot B_a.
  \]
  Then there is a projective submeasurement $P=\{P_a\}$ with
  \[
    A_a \ot I \approx_{84\zeta^{1/4}} P_a \ot I.
  \]
\end{lemma}
\begin{proof}
  Follow the Kempe--Vidick orthogonalization argument~\cite{KV11}. First the consistency with $B$ is converted into an ``almost projective'' estimate for $A$, namely that $\sum_a \bra{\psi} (A_a-A_a^2) \ot I \ket{\psi}$ is small. Spectrally truncate each $A_a$ to a projector $R_a$, then discard the small singular directions to obtain lower-rank projectors $Q_a$ whose sum $Q=\sum_a Q_a$ is still close to a projector but may not equal the identity. The source then rewrites $Q_a$ as $X^\dagger T_a X$ for a genuine projective measurement $T$, replaces $X$ by the unitary part $\widehat X$ of its singular value decomposition, and defines $P_a=\widehat X^\dagger T_a \widehat X$. The main bookkeeping burden is to compare successively $A_a$, $R_a$, $Q_a$, and $P_a$ while controlling the completeness defects of $Q$ and $\sqrt Q$; in Lean these appear as several explicit helper lemmas rather than one black-box step.
\end{proof}

\begin{theorem}[Orthogonalization lemma]\label{thm:orthonormalization}
  \lean{MIPStarRE.Paper2009LDT.Section5MakingMeasurementsProjective.orthonormalization}
  \uses{def:strong-self-consistency, def:approx_delta}
  Let $A=\{A_a\}$ be a submeasurement on a permutation-invariant state $\ket{\psi}$ such that
  \[
    \sum_a \bra{\psi} A_a \ot A_a \ket{\psi}
    \ge
    \sum_a \bra{\psi} A_a \ot I \ket{\psi} - \zeta.
  \]
  Then there is a projective submeasurement $P=\{P_a\}$ with
  \[
    A_a \ot I \approx_{100\zeta^{1/4}} P_a \ot I.
  \]
\end{theorem}
\begin{proof}
  \uses{prop:other-two-notions-of-self-consistency, def:measurement-completion, lem:orthonormalization-main-lemma}
  Complete $A$ by adding the missing mass as an extra outcome. The completed measurement is $2\zeta$-self-consistent by Lemma~\ref{prop:other-two-notions-of-self-consistency}, so Lemma~\ref{lem:orthonormalization-main-lemma} applies. Finally discard the extra outcome.
\end{proof}
